
\newpage
\subsection{滤过 $A_\infty$ 双模与凝聚层复形}
\label{subsec:5-4-3}

下面解释滤过 $A_\infty$ 双模如何在形式概形
\[
\widehat{\mathcal M}(C,\mathfrak m)
\times_{\Lambda_{0,\mathrm{nov}}}
\widehat{\mathcal M}(C',\mathfrak m')
\]
上定义凝聚层导出范畴中的一个对象。
\translatornote{原书在本小节若干处
把第二系数代数 $C'$ 的结构写成
$\mathfrak m$ 而非 $\mathfrak m'$；
这里依第 5.4.1 小节已固定的
带撇记号统一，见疑误 A-0320。}

先固定记号。
设 $D$ 是以 $\Lambda_{0,\mathrm{nov}}$
为系数的 $(C,C')$-滤过
$A_\infty$ 双模，
并仍假设 $C,C',D$ 都有隙。
本小节为简化记号，
用 $(C,C')$ 代替 $(C_1,C_0)$。
\translatornote{原书实体页 346
把 $C$ 的奇基数目写作 $\ell$、
偶基数目写作 $m$；
但定义 5.4.7--5.4.8
已固定相反的指标分配，
且本页稍后的形式环也照此书写。
以下统一为 $m$ 个
$\mathbf v_i,s_i$ 与
$\ell$ 个 $\mathbf w_j,h_j$，
带撇侧同理；见疑误 A-0318。}
取 $C^{\mathrm{odd}}$ 的基
$\mathbf v_1,\ldots,\mathbf v_m$、
$C^{\mathrm{even}}$ 的基
$\mathbf w_1,\ldots,\mathbf w_\ell$，
相应形式坐标分别为
$s_1,\ldots,s_m$
和 $h_1,\ldots,h_\ell$。
对 $C'$ 使用带撇的对应记号：
$\mathbf v'_1,\ldots,\mathbf v'_{m'}$、
$\mathbf w'_1,\ldots,\mathbf w'_{\ell'}$
及 $s'_1,\ldots,s'_{m'}$、
$h'_1,\ldots,h'_{\ell'}$。
令
\[
\begin{aligned}
\mathcal R_C
&=\Lambda_{0,\mathrm{nov}}
[[s_1,\ldots,s_m;h_1,\ldots,h_\ell]],\\
\mathcal R_{C'}
&=\Lambda_{0,\mathrm{nov}}
[[s'_1,\ldots,s'_{m'};h'_1,\ldots,h'_{\ell'}]],\\
z
&=s_1\mathbf v_1+\cdots+s_m\mathbf v_m
+h_1\mathbf w_1+\cdots+h_\ell\mathbf w_\ell,\\
z'
&=s'_1\mathbf v'_1+\cdots+s'_{m'}\mathbf v'_{m'}
+h'_1\mathbf w'_1+\cdots+h'_{\ell'}\mathbf w'_{\ell'}.
\end{aligned}
\]
设 $D\cong\Lambda_{0,\mathrm{nov}}^{\oplus d}$，
并取基 $\mathbf p_1,\ldots,\mathbf p_d$。
在拉格朗日交弗洛尔上同调的几何情形中，
这些基元对应于一对拉格朗日子流形的交点，
或 $L^{(1)}\cap L^{(0)}$
的上同调群的一组基。
以 $\pi_i\colon D\to\Lambda_{0,\mathrm{nov}}$
表示取 $\mathbf p_i$ 坐标的投影。

\begin{definition}{5.4.22}\label{def:5-4-22}
定义
\begin{equation}
\begin{aligned}
Q_{i;j}
(&s_1,\ldots,s_m;h_1,\ldots,h_\ell;
s'_1,\ldots,s'_{m'};h'_1,\ldots,h'_{\ell'})
\\
&=(\pi_j\circ\mathfrak n)
\bigl(\exp z,\mathbf p_i,\exp z'\bigr)
\in\mathcal R_C
\widehat\otimes_{\Lambda_{0,\mathrm{nov}}}
\mathcal R_{C'}.
\end{aligned}
\tag{5.4.23}\label{eq:5-4-23}
\end{equation}
\end{definition}

(5.4.23) 右端的理解方式与 (5.4.10) 相同。
映射
\[
\mathbf p_i
\longmapsto
\sum_{j=1}^{d}
Q_{i;j}\cdot\mathbf p_j
\]
定义同态
\[
\mathfrak Q_0\colon
D\longrightarrow
\mathcal R_C\widehat\otimes
D\widehat\otimes\mathcal R_{C'}.
\]
而且，$\mathfrak Q_0$ 诱导双模同态
\[
\mathfrak Q_1\colon
\mathcal R_C\widehat\otimes D
\widehat\otimes\mathcal R_{C'}
\longrightarrow
\mathcal R_C\widehat\otimes D
\widehat\otimes\mathcal R_{C'},
\]
其作用为
\[
\mathfrak Q_1(x,y,z)
=(-1)^{\deg' x}\,
x\otimes\mathfrak Q_0(y)\otimes z.
\]
这里及以下的完备张量积
均以 $\Lambda_{0,\mathrm{nov}}$ 为底。

再记
\[
\begin{aligned}
\mathcal R_C^{\mathrm{bos}}
&=\Lambda_{0,\mathrm{nov}}[[s_1,\ldots,s_m]],\\
\mathcal R_{C'}^{\mathrm{bos}}
&=\Lambda_{0,\mathrm{nov}}[[s'_1,\ldots,s'_{m'}]].
\end{aligned}
\]
定义
\[
\begin{aligned}
\mathfrak M(D,\mathfrak n)
&=
\frac{\mathcal R_C}
{(P_1,\ldots,P_{m+\ell})}
\widehat\otimes D\widehat\otimes
\frac{\mathcal R_{C'}}
{(P'_1,\ldots,P'_{m'+\ell'})},\\[2mm]
\mathfrak M(D,\mathfrak n)^{\mathrm{bos}}
&=
\frac{\mathcal R_C^{\mathrm{bos}}}
{(P_{m+1}^{\mathrm{bos}},
\ldots,P_{m+\ell}^{\mathrm{bos}})}
\widehat\otimes D\widehat\otimes
\frac{\mathcal R_{C'}^{\mathrm{bos}}}
{(P_{m'+1}^{\prime\mathrm{bos}},
\ldots,P_{m'+\ell'}^{\prime\mathrm{bos}})}.
\end{aligned}
\]
这里的 $P_i$ 和 $P'_i$
分别是 $C$ 和 $C'$ 的超仓西映射。\translatornote{原书实体页 346
在上面两个右侧商理想的末项
分别漏写 $P'_{m'+\ell'}$
及 $P_{m'+\ell'}^{\prime\mathrm{bos}}$
的撇号；
右侧是 $C'$ 的形式环，
故按同一列首项和式后定义补回。
见疑误 A-0319。}

显然，$\mathfrak M(D,\mathfrak n)$
是
$(\mathfrak M(C,\mathfrak m),
\mathfrak M(C',\mathfrak m'))$
双模，而 $\mathfrak Q_1$
诱导双模同态
\[
\mathfrak Q\colon
\mathfrak M(D,\mathfrak n)
\longrightarrow
\mathfrak M(D,\mathfrak n).
\]

\begin{proposition}{5.4.24}\label{prop:5-4-24}
$\mathfrak Q^2=0$。
\end{proposition}

\begin{proof}
记
\[
\boldsymbol s\boldsymbol v
=s_1\mathbf v_1+\cdots+s_m\mathbf v_m
+h_1\mathbf w_1+\cdots+h_\ell\mathbf w_\ell
=z,
\]
并对带撇对象同样记
$\boldsymbol s'\boldsymbol v'=z'$。
由 $\mathfrak n$ 的
$A_\infty$ 双余模结构，得到
\[
\begin{aligned}
0={}&
\mathfrak n\bigl(
\exp(\rho z),
\mathfrak n(\exp z,\mathbf p_i,\exp z'),
\exp z'
\bigr)\\
&+
\mathfrak n\bigl(
\widehat d_C(\exp z),
\mathbf p_i,\exp z'
\bigr)\\
&+
(-1)^{\deg\mathbf p_i}
\mathfrak n\bigl(
\exp(\rho z),
\mathbf p_i,
\widehat d_{C'}(\exp z')
\bigr).
\end{aligned}
\]
这里 $\rho$ 由
$\rho(\mathbf v_i)=\mathbf v_i$、
$\rho(\mathbf w_i)=-\mathbf w_i$
定义；
$\widehat d_C$、
$\widehat d_{C'}$
分别表示两侧代数的余导子。
按超仓西映射的定义，
第二、第三项模由两侧
超仓西映射生成的理想为零。
第一项正给出
$\mathfrak Q^2(\mathbf p_i)$，
故命题得证。\translatornote{原书实体页 347
把第一项说成
$\mathfrak Q^2(\boldsymbol s\boldsymbol v)$；
但 $\boldsymbol s\boldsymbol v$
是代数输入，
而 $\mathfrak Q$ 作用于以 $D$
为中间因子的双模。
此处应在基元 $\mathbf p_i$
上验证，见疑误 A-0323。}
\end{proof}

由此，
$(\mathfrak M(D,\mathfrak n),\mathfrak Q)$
成为下列两个微分分次代数上的微分双模：
\[
\mathfrak M(C,\mathfrak m),
\qquad
\mathfrak M(C',\mathfrak m').
\]
特别地，它给出
\[
\mathfrak M(C,\mathfrak m)^{\mathrm{bos}}
\otimes_{\Lambda_{0,\mathrm{nov}}}
\mathfrak M(C',\mathfrak m')^{\mathrm{bos}}
\]
上的带微分模。
因而，在形式概形
\liftIndex{I213}
\[
\operatorname{Spf}
\bigl(\mathfrak M(C,\mathfrak m)^{\mathrm{bos}}\bigr)
\times_{\Lambda_{0,\mathrm{nov}}}
\operatorname{Spf}
\bigl(\mathfrak M(C',\mathfrak m')^{\mathrm{bos}}\bigr)
\]
上，它定义一个凝聚层复形
$\mathfrak M(D,\mathfrak n)^{\sim}$
（见 \textbf{[GrDi60]}）。
\liftIndex{I212}
若 $D$ 有限生成，
这个复形是凝聚的。

令
$b\in\widehat{\mathcal M}(C,\mathfrak m)$、
$b'\in\widehat{\mathcal M}(C',\mathfrak m')$。
由命题 5.4.15，
可将它们分别看作连续同态
\[
\begin{aligned}
i_b&\colon
\mathfrak M(C,\mathfrak m)
\longrightarrow\Lambda_{0,\mathrm{nov}},\\
i_{b'}&\colon
\mathfrak M(C',\mathfrak m')
\longrightarrow\Lambda_{0,\mathrm{nov}}.
\end{aligned}
\]
它们在 $\Lambda_{0,\mathrm{nov}}$
上给出相应的双模结构；
将第一个点值右模记为
$(\Lambda_{0,\mathrm{nov}})_b$，
第二个点值左模记为
\({}_{b'}(\Lambda_{0,\mathrm{nov}})\)。
\translatornote{原书实体页 347
把这两个点值因子依次称为
左模、右模，
与紧接的张量顺序相反；
又把下式称为层的“茎”，
而该张量基变换实际给出点处纤维。
见疑误 A-0321、A-0324。}
按定义（见 \textbf{[GrDi60]}），
层 $\mathfrak M(D,\mathfrak n)^{\sim}$
在 $(b,b')$ 处的纤维是
\[
\begin{aligned}
\mathfrak M(D,\mathfrak n)^{\sim}_{b,b'}
={}&
(\Lambda_{0,\mathrm{nov}})_b
\otimes_{\mathfrak M(C,\mathfrak m)}
\mathfrak M(D,\mathfrak n)\\
&{}\otimes_{\mathfrak M(C',\mathfrak m')}
{}_{b'}(\Lambda_{0,\mathrm{nov}}).
\end{aligned}
\]
作为 $\Lambda_{0,\mathrm{nov}}$-模，
它同构于 $D$；
但 $\mathfrak Q$ 在其上诱导的微分
依赖于 $b,b'$。
以下引理立即由定义得到。

\begin{lemma}{5.4.25}\label{lem:5-4-25}
$\mathfrak Q$ 在
$\mathfrak M(D,\mathfrak n)^{\sim}_{b,b'}$
上诱导的微分正是
$\delta_{b,b'}$，其中
\[
\delta_{b,b'}(\mathbf p)
=\mathfrak n
\bigl(e^b,\mathbf p,e^{b'}\bigr).
\]
\end{lemma}

还有如下结果。
\begin{lemma}{5.4.26}\label{lem:5-4-26}
若两个 $(C,C')$-滤过
$A_\infty$ 双模 $D,D'$
按定义 5.4.3(3) 同构，
则
$\mathfrak M(D,\mathfrak n)^{\sim}$
与
$\mathfrak M(D',\mathfrak n')^{\sim}$
作为凝聚层复形同构。
\end{lemma}

\begin{proof}
令
$\varphi\colon D\to D'$、
$\varphi^{-1}\colon D'\to D$
为互逆的滤过 $A_\infty$
双模同态，
均覆盖
$(\mathrm{id}_C,\mathrm{id}_{C'})$。\translatornote{原书实体页 348
在该引理的证明中，
把覆盖任意 $(f,g)$ 的同态
用于定义 5.4.3(3)，
又把诱导映射的目标误写回
$\mathfrak M(D,\mathfrak n)^\sim$。
该定义只涉及覆盖恒等的双模同构；
下文按这一陈述修正目标与基底，
见疑误 A-0322。}
对每一对
\[
(b,b')\in
\widehat{\mathcal M}(C,\mathfrak m)
\times
\widehat{\mathcal M}(C',\mathfrak m'),
\]
$\varphi$ 自然诱导点处纤维同构
\[
\mathfrak M(D,\mathfrak n)^{\sim}_{b,b'}
\longrightarrow
\mathfrak M(D',\mathfrak n')^{\sim}_{b,b'}.
\]
这些同构的自然性
给出凝聚层复形的同构。
\end{proof}

在几何情形中，这蕴涵：
\begin{corollary}{5.4.27}\label{cor:5-4-27}
设 $(L^{(1)},L^{(0)})$
是由无障碍拉格朗日子流形
组成的相对 spin 对。
则存在凝聚层导出范畴中的对象
\[
\bigl(
\mathfrak M(L^{(1)},L^{(0)})^{\sim},
\mathfrak Q
\bigr),
\]
它在每个
$\Lambda^+_{0,\mathrm{nov}}$-值
（刚性）点处给出弗洛尔上同调。
这个对象不依赖于相容殆复结构的选择，
也不依赖于定义过程中涉及的其他选择。
\end{corollary}
